% AI-produced by GPT-6 Astra (OpenAI, ChatGPT Pro mode).
% Recovered from the actual CP14 draft; root and independent audits passed in-session.
% Promotion edits restored on 2026-11-20; original draft retained unchanged.
\Needspace{10\baselineskip}
\subsection{A.142. A component obstruction in real base dimension three}

\begin{saresult}{A.142}{Proved under stated hypotheses}{Failure of PL--TOP surjectivity on concordance components over the projective three-space}
Let $B$ be a closed, connected, oriented real PL three-manifold whose
fundamental group is finite cyclic of even order. There is a topological
homeomorphism
\[
 F:B\times I\longrightarrow B\times I,\qquad
 F|_{B\times\{0,1\}}=1,
                                                        \tag{G.1}
\]
whose component in the ordinary free-top concordance group
$P_{\mathrm{TOP}}(B)$ contains no PL concordance. In particular,
\[
 \pi_0P_{\mathrm{PL}}(\mathbb {RP}^3)\longrightarrow
 \pi_0P_{\mathrm{TOP}}(\mathbb {RP}^3)
 \quad\hbox{is not surjective}.                         \tag{G.2}
\]
For every fixed o-minimal expansion of the real field, the natural map
$\pi_0P_{\mathrm{def}}(B)\to\pi_0P_{\mathrm{TOP}}(B)$ is likewise
not surjective. This includes the semialgebraic category.
Here PL families are jointly PL, TOP families are jointly continuous,
and every concordance fixes the bottom pointwise. The obstruction is a
nonzero element of
\[
 H^3(B\times I,B\times\{0,1\};\mathbb Z/2)
       \cong H^2(B;\mathbb Z/2).
                                                        \tag{G.3}
\]
No assertion that vanishing of this obstruction suffices for an
unstabilized PL isotopy is made.
\end{saresult}
\begin{proof}
\textbf{The relative PL obstruction and the category convention.}
Put $X=B\times I$. Choose a compatible smoothing of $B$, and its product
smoothing on $X$. Existence of a single compatible smoothing in these
dimensions is the classical low-dimensional PL smoothing theorem; see
\cite[Essay I, p.~39, final paragraph]{KirbySiebenmann1977}. No
comparison of four-dimensional smooth and PL isotopy groups is used.

For a homeomorphism $f:X\to X$ which fixes $\partial X$ pointwise, form
the five-dimensional mapping cylinder
\[
 M_f=(X\times I)\mathbin{\cup_f}X.
                                                        \tag{G.4}
\]
The given PL structure of $X$ determines a PL reduction of the stable
tangent microbundle on both ends of $M_f$, and the identity restriction
of $f$ makes these agree along the side boundary. Let
\[
 c_{\mathrm{PL}}(f)=
 \sigma^{-1}k(M_f,\partial M_f)
       \in H^3(X,\partial X;\mathbb Z/2),
                                                        \tag{G.5}
\]
where $k$ is the relative universal Kirby--Siebenmann obstruction to
extending this boundary PL reduction, and $\sigma$ is the relative
interval-suspension isomorphism. This is a formal tangent-bundle
obstruction. Its definition requires no theorem classifying actual PL
structures or isotopies on a four-manifold.

The class in (G.5) agrees with Galvin's Casson--Sullivan invariant for
the chosen compatible smoothing. Indeed, his Definition~2.11 and
equation~(2.1), printed pp.~2503--2504, define the latter by this same
relative mapping-cylinder obstruction and suspension. The natural map
from smooth reductions to PL reductions preserves the universal
Kirby--Siebenmann class; the low-dimensional identification used here is
explicit in his Remark~2.5 and \S2.2. Thus, for the boundary-identity maps
under consideration,
\[
             c_{\mathrm{PL}}(f)=\operatorname{cs}(f).
                                                        \tag{G.6}
\]
This identifies the classes, rather than asserting an equivalence of
smooth and PL isotopies.

Two elementary properties will be needed. First, $c_{\mathrm{PL}}$ is
constant under TOP isotopies fixed on $\partial X$: an isotopy provides
the corresponding relative homotopy of the tangent-microbundle
identifications and hence of the relative obstruction classes. One can
equivalently use the mapping-cylinder proof of
\cite[Proposition~2.22, printed pp.~2507--2508]{Galvin2026}, with PL
reductions throughout. Boundary collars may first be made common for
the whole compact parameter interval by the compression construction
of A.135. Second, if $f$ itself is PL, its mapping cylinder has a PL
structure extending exactly these boundary structures, so
\[
                 f\text{ PL}\quad\Longrightarrow\quad
                 c_{\mathrm{PL}}(f)=0.
                                                        \tag{G.7}
\]
The second fact applies to an actual PL map even when it has not been
identified with a smooth map in fixed coordinates.

\textbf{The compact relative realization, including the side boundary.}
We spell out the boundary version of Galvin's construction. This also
resolves a wording discrepancy in the source: Theorem~1.7 and
Proposition~4.13 allow compact manifolds, whereas Definition~4.11 and
the opening of \S4.4 say ``closed''. The argument below uses the relative
normal invariants and the relative product theorem of \S4.2 directly.

For a compact oriented smooth four-manifold $X$ with finite cyclic
fundamental group, the simple surgery obstruction group is zero:
\[
                 L^s_5(\mathbb Z[\pi_1X])=0.
                                                        \tag{G.8}
\]
This is the finite-cyclic input stated in
\cite[\S4.4, printed p.~2528]{Galvin2026}, citing Bak's calculation.
Finite cyclic groups are good for four-dimensional topological surgery
and the relative topological $s$-cobordism theorem. The simple relative
surgery sequence therefore realizes every normal invariant of
$(X\times I,\partial(X\times I))$ by a simple structure. In the exact
normal-invariant coordinates of Galvin's Lemmas~4.8--4.9, printed
p.~2526, these groups and their PL obstruction map are
\[
 \begin{split}
 [(X\times I,\partial),G/\mathrm{TOP}]&\cong
 H^2(X\times I,\partial;\mathbb Z/2)
       \oplus H^4(X\times I,\partial;\mathbb Z),\\
 \xi_*(x,y)&=\operatorname{red}_2(y)
       \in H^4(X\times I,\partial;\mathbb Z/2).
 \end{split}                                            \tag{G.9}
\]
For completeness, the potential additional term is $\operatorname{Sq}^2x$.
Here $X\times I/\partial(X\times I)$ is the reduced suspension of
$X/\partial X$, so $x$ is suspended from a degree-one relative class.
Steenrod squares commute with suspension, and $\operatorname{Sq}^2$ on
a degree-one class is zero by the unstable axiom. The final inequality
in the printed proof of Lemma~4.9 is reversed; the correct axiom is
$\operatorname{Sq}^i u=0$ when $i>\deg u$. The displayed formula (G.9)
and its relative proof are unaffected.

The coefficient reduction in degree four is onto. In fact the relevant
Bockstein segment is
\[
 H^4(X\times I,\partial;\mathbb Z)\longrightarrow
 H^4(X\times I,\partial;\mathbb Z/2)\longrightarrow
 H^5(X\times I,\partial;\mathbb Z)
       \xrightarrow{\;2\;}
 H^5(X\times I,\partial;\mathbb Z),
                                                        \tag{G.10}
\]
and the last group is $\mathbb Z$ by relative orientation. Given any
$\eta\in H^3(X,\partial X;\mathbb Z/2)$, choose $y$ reducing to
$\sigma\eta$, and use the normal invariant $(0,y)$. Its surgery
obstruction is zero by (G.8).

We now check that its resulting self-homeomorphism is the identity on
$\partial X$, not merely a map preserving that boundary. Represent
the lifted simple structure by
\[
 \phi:(N;\partial_+N,\partial_0N,\partial_-N)
   \longrightarrow(X\times I;X\times0,\partial X\times I,X\times1),
                                                        \tag{G.11}
\]
where $\phi$ is a simple homotopy equivalence and its full boundary
restriction is the prescribed homeomorphism. Use this restriction to
identify the side $\partial_0N$ with $D\times I$, where
$D=\partial(\partial_+N)$, compatibly with its bottom identification.
Galvin's Proposition~4.10, printed pp.~2526--2527, applies the relative
topological $s$-cobordism theorem to obtain
\[
 p:N\longrightarrow\partial_+N\times I
 \quad\hbox{equal to the identity on }
       \partial_+N\cup(D\times I).
                                                        \tag{G.12}
\]
Write $p_-$ for its top restriction and $\phi_\pm$ for the two boundary
labelings by $X$. The self-homeomorphism in the source's commuting
diagram is
\[
                  f=\phi_+p_-\phi_-^{-1}:X\longrightarrow X.
                                                        \tag{G.13}
\]
On $\partial X$, the two labelings in (G.13) agree after $p_-$ precisely
because (G.12) fixes the labeled side product. Consequently
$f|_{\partial X}=1$. This is the only extra boundary observation needed
in passing from the displayed source construction to (G.1).
The induced identification of $N$ with $M_f$ preserves the specified
boundary reductions. Lemma~4.12, printed p.~2527, or directly (G.5),
therefore gives
\[
                   c_{\mathrm{PL}}(f)=\eta.
                                                        \tag{G.14}
\]
This proves relative realization without applying the word ``closed''
in Definition~4.11 to a manifold with boundary.

\textbf{Why freeing the top cannot remove a nonzero class.}
For every compact PL three-manifold $B$, the natural map of actual
joint-family automorphism groups
\[
 \mathcal H_{\mathrm{PL}}(B)\longrightarrow
 \operatorname{Sing}\operatorname{Homeo}(B)
                                                        \tag{G.15}
\]
is a weak equivalence by (A.134.2). Its proof uses Hamstrom's
Theorem~3.1.2 with the joint PL convention of her \S2.1, together with
the explicit all-components argument in A.134. It does not use the
block convention from the earlier quasi-simplicial comparison.

An immediate, useful consequence is the following relative path
statement. If $a_t$ is a TOP path from the identity to a PL map $a_1$,
there is a jointly PL path $b_t$ with the same endpoints and a TOP
homotopy of paths from $a_t$ to $b_t$ fixed at both endpoints. Indeed,
injectivity on $\pi_0$ supplies some PL path $b^0$. Surjectivity on
$\pi_1$ supplies a PL loop representing $a*(b^0)^{-1}$. Concatenating
that loop with $b^0$ gives $b$. These are finite PL concatenations;
their adjoints on $B\times I$ are jointly PL.

Suppose the boundary-identity realization $F$ of a nonzero $\eta$ were
joined to a PL concordance $G$ by a path $H_t$ in the ordinary
free-top TOP group, with $H_0=F$ and $H_1=G$. Put
$a_t=H_t|_{B\times1}$, so $a_0=1$ and $a_1$ is PL. Define
\[
 T_t(x,s)=(a_{ts}(x),s),\qquad Q_t=T_t^{-1}H_t.
                                                        \tag{G.16}
\]
Every $T_t$ is a homeomorphism with continuous inverse, jointly in all
parameters, because $B$ is compact. The path $Q_t$ fixes both ends
pointwise, starts at $F$, and ends at $(T_1)^{-1}G$.
Replace $a$ by the jointly PL path $b$ just obtained. The homotopy of
paths relative endpoints gives a TOP isotopy, fixed on the entire
boundary of $X$, from $T_1$ to
\[
                       T^b(x,s)=(b_s(x),s).
                                                        \tag{G.17}
\]
The latter is a PL homeomorphism. Hence $(T_1)^{-1}G$ is TOP-isotopic
relative to the boundary to the PL map $(T^b)^{-1}G$, which fixes that
boundary pointwise. Combining these paths and (G.7) yields
$c_{\mathrm{PL}}(F)=0$, contrary to (G.14). This argument only uses the
necessity of vanishing for a PL map and isotopy invariance.

\textbf{The specified base and its finite certificate.}
For a closed oriented three-manifold with finite cyclic fundamental
group of even order, Poincare duality gives
\[
 H^2(B;\mathbb Z/2)\cong H_1(B;\mathbb Z/2)\cong\mathbb Z/2.
                                                        \tag{G.18}
\]
Choose the nonzero class and apply the preceding construction. In
particular, $B=\mathbb {RP}^3$ is oriented, has fundamental group
$\mathbb Z/2$, and has the required nonzero class.

The definable conclusion follows with its natural real comparison
map, rather than through an abstract identification of groups. Realize
$B$ as a finite real polyhedron, which is semialgebraic. By A.113 every
definable concordance is definably isotopic to a PL concordance, relative
to the bottom. Such an isotopy is an ordinary TOP isotopy. Thus the
images of the PL and definable component maps in $\pi_0P_{\mathrm{TOP}}(B)$
coincide, and the component of $F$ belongs to neither image. This
argument applies separately to every fixed o-minimal expansion of
$\mathbb R$; it asserts no comparison map from a non-Archimedean field.

For an explicit finite PL base take the barycentric subdivision of
the boundary of the four-dimensional crosspolytope and quotient by
the antipodal simplicial involution. The quotient realizes the usual
two-fold covering $S^3\to\mathbb {RP}^3$. It is a PL manifold: the
antipodal action has disjoint local stars after subdivision and the
quotient map is locally a PL isomorphism. The universal covering
$S^3$ is simply connected, proving $\pi_1=\mathbb Z/2$ independently
of a homology calculation.

The exact finite certificate
\path{code/check_projective_three_base.py}, with output
\path{logs/check_projective_three_base.json}, lists all 192
tetrahedra and checks the face vector
$(40,232,384,192)$, the mod-two Betti numbers $(1,1,1,1)$,
80 local maximal-star bijections for the two-fold cover, and an
edge-path presentation reduced by elementary Tietze eliminations to
one generator with square relations. The parent's recorded run
passed 2634 checks in 0.08755117800319567 seconds; the author independently
read and reran it, obtaining 2634 passes in 0.10303335600474384 seconds,
recorded in
\path{agents/singularities_duality/logs/check_projective_three_base_author.json}.
This is a finite
certificate for the base and the nonzero obstruction group; the
existence and exclusion of $F$ are established by the proof above.
\end{proof}

\textbf{Literature concordance.}
Galvin's 2024 thesis, \emph{Non-smoothable homeomorphisms of
4-manifolds}, Theorem~9.0.1 and Corollary~9.0.2, printed p.~109,
already establish non-surjectivity on the first homotopy groups of
certain block-automorphism comparison maps; Remark~9.0.3, p.~110,
includes even-order lens spaces. The proof in \S9.3, printed
pp.~115--117, uses boundary-identity homeomorphisms of the product
with an interval and their Casson--Sullivan invariants.
The block convention is defined in \S9.1, pp.~110--112.
Our statement above is about the ordinary free-top concordance
component group and its actual PL and definable images. The endpoint
trace and joint-family comparison in (G.15)--(G.17) supply that
passage explicitly. No identification of block loops with ordinary
loops, and no zero-invariant sufficiency assertion, is inferred.
See \cite[Theorem~9.0.1, Corollary~9.0.2, and \S9.3]{GalvinThesis2024}.
